% GENERATED FILE. Edit canonical lessons, then run npm run generate:books.
% canonical-source-hash: fnv1a64:6dfa7ee902ac86a2
% canonical-lessons: FR-C358-syndicat, FR-C358-greve, FR-C358-greviste, FR-C358-salarie, FR-C358-heures-supplementaires

\chapter{La grève}
\label{ch:fr-la-greve}
\hlchaptermodality{358}

\begin{chapteropening}
\textbf{By the end of this chapter:} \emph{I can talk about workers' rights: the union, a strike, strikers, employees and overtime.}

Complete the last lesson of chapter 358: I can talk about workers' rights: the union, a strike, strikers, employees and overtime.
\end{chapteropening}

\section[le syndicat]{le syndicat — the trade union}

\label{lesson:FR-C358-syndicat}



\begin{quote}
Before the new one: say the French for the apprenticeship, then the French for the employer.
\end{quote}

\subsection*{You'll want to know: le syndicat}
\textbf{le syndicat} — \textquotedblleft{}the trade union\textquotedblright{}.

\begin{quote}
Mon père travaille pour un syndicat.
\end{quote}

\begin{quote}
\emph{My father works for a trade union.}
\end{quote}

\emph{Le syndicat d'initiative} is an old name for a tourist office.

\begin{cousinweb}
\emph{syndicat} is built from \emph{syndic} \textquotedblleft{}representative\textquotedblright{} and \emph{-at} (EtymDB). \emph{syndic} comes from Late Latin \emph{syndicus} \textquotedblleft{}delegate of a body\textquotedblright{}, from Greek \emph{sýndikos} \textquotedblleft{}advocate\textquotedblright{}.
\end{cousinweb}

\begin{grammarlens}[title={What you've built}]
One more word for telling, explaining and describing at B1.
\end{grammarlens}

\subsection*{Your turn}
\begin{itemize}
\raggedright
  \item \emph{Say it:} \emph{le syndicat}
  \item \emph{Say it:} \emph{le syndicat}, once more
  \item \emph{Say it:} a sentence of your own with \emph{le syndicat}
  \item \emph{Recall:} say the French for the apprenticeship, then the French for the employer, then say \emph{le syndicat} again
\end{itemize}

\begin{tcolorbox}[breakable,colback=teal!4,colframe=teal!35!black,title={Before you move on}]
What does le syndicat mean? (\textquotedblleft{}The trade union\textquotedblright{}.) Say it once more.
\end{tcolorbox}
\section[la grève]{la grève — the strike}

\label{lesson:FR-C358-greve}



\begin{quote}
Before the new one: say the French for the employer, then the French for the trade union.
\end{quote}

\subsection*{You'll want to know: la grève}
\textbf{la grève} — \textquotedblleft{}the strike\textquotedblright{}.

\begin{quote}
Les trains ne roulent pas : il y a une grève.
\end{quote}

\begin{quote}
\emph{The trains are not running: there is a strike.}
\end{quote}

\emph{Faire grève}, \emph{être en grève} — to be on strike. \emph{La grève} is also a sandy shore.

\begin{cousinweb}
\emph{grève} \textquotedblleft{}strike\textquotedblright{} is built on \emph{grève} \textquotedblleft{}sandy shore\textquotedblright{}: the usual explanation is a riverside square of that name where people looking for work gathered. The origin of \emph{grève} \textquotedblleft{}shore\textquotedblright{} is not settled.
\end{cousinweb}

\begin{grammarlens}[title={What you've built}]
One more word for telling, explaining and describing at B1.
\end{grammarlens}

\subsection*{Your turn}
\begin{itemize}
\raggedright
  \item \emph{Say it:} \emph{la grève}
  \item \emph{Say it:} \emph{la grève}, once more
  \item \emph{Say it:} a sentence of your own with \emph{la grève}
  \item \emph{Recall:} say the French for the employer, then the French for the trade union, then say \emph{la grève} again
\end{itemize}

\begin{tcolorbox}[breakable,colback=teal!4,colframe=teal!35!black,title={Before you move on}]
What does la grève mean? (\textquotedblleft{}The strike\textquotedblright{}.) Say it once more.
\end{tcolorbox}
\section[le gréviste]{le gréviste — the striker}

\label{lesson:FR-C358-greviste}



\begin{quote}
Before the new one: say the French for the trade union, then the French for the strike.
\end{quote}

\subsection*{You'll want to know: le gréviste}
\textbf{le gréviste} — \textquotedblleft{}the striker\textquotedblright{}.

\begin{quote}
Les grévistes demandent de meilleurs salaires.
\end{quote}

\begin{quote}
\emph{The strikers are asking for better wages.}
\end{quote}

The same form for a man or a woman.

\begin{cousinweb}
\emph{gréviste} is built from \emph{grève} \textquotedblleft{}strike\textquotedblright{} and \emph{-iste}. \emph{grève} first meant a sandy shore; its origin is not settled.
\end{cousinweb}

\begin{grammarlens}[title={What you've built}]
One more word for telling, explaining and describing at B1.
\end{grammarlens}

\subsection*{Your turn}
\begin{itemize}
\raggedright
  \item \emph{Say it:} \emph{le gréviste}
  \item \emph{Say it:} \emph{le gréviste}, once more
  \item \emph{Say it:} a sentence of your own with \emph{le gréviste}
  \item \emph{Recall:} say the French for the trade union, then the French for the strike, then say \emph{le gréviste} again
\end{itemize}

\begin{tcolorbox}[breakable,colback=teal!4,colframe=teal!35!black,title={Before you move on}]
What does le gréviste mean? (\textquotedblleft{}The striker\textquotedblright{}.) Say it once more.
\end{tcolorbox}
\section[le salarié]{le salarié — the employee, the wage earner}

\label{lesson:FR-C358-salarie}



\begin{quote}
Before the new one: say the French for the strike, then the French for the striker.
\end{quote}

\subsection*{You'll want to know: le salarié}
\textbf{le salarié} — \textquotedblleft{}the employee, the wage earner\textquotedblright{}.

\begin{quote}
Cette entreprise a deux cents salariés.
\end{quote}

\begin{quote}
\emph{This company has two hundred employees.}
\end{quote}

From \emph{salaire}. Feminine \emph{salariée}.

\begin{cousinweb}
\emph{salarié} is built on \emph{salaire} \textquotedblleft{}wages\textquotedblright{}, from Latin \emph{salārium} \textquotedblleft{}salt money, an allowance\textquotedblright{} (Lewis \& Short), from \emph{sāl} \textquotedblleft{}salt\textquotedblright{}.
\end{cousinweb}

\begin{grammarlens}[title={What you've built}]
One more word for telling, explaining and describing at B1.
\end{grammarlens}

\subsection*{Your turn}
\begin{itemize}
\raggedright
  \item \emph{Say it:} \emph{le salarié}
  \item \emph{Say it:} \emph{le salarié}, once more
  \item \emph{Say it:} a sentence of your own with \emph{le salarié}
  \item \emph{Recall:} say the French for the strike, then the French for the striker, then say \emph{le salarié} again
\end{itemize}

\begin{tcolorbox}[breakable,colback=teal!4,colframe=teal!35!black,title={Before you move on}]
What does le salarié mean? (\textquotedblleft{}The employee, the wage earner\textquotedblright{}.) Say it once more.
\end{tcolorbox}
\section[les heures supplémentaires]{les heures supplémentaires — the overtime}

\label{lesson:FR-C358-heures-supplementaires}



\begin{quote}
Before the new one: say the French for the striker, then the French for the employee.
\end{quote}

\subsection*{You'll want to know: les heures supplémentaires}
\textbf{les heures supplémentaires} — \textquotedblleft{}the overtime\textquotedblright{}.

\begin{quote}
Ce mois-ci, j'ai fait vingt heures supplémentaires.
\end{quote}

\begin{quote}
\emph{This month, I did twenty hours of overtime.}
\end{quote}

Often shortened to \emph{les heures sup}.

\begin{cousinweb}
\emph{supplémentaire} is built on \emph{supplément}, from Latin \emph{supplēmentum} \textquotedblleft{}a filling up, a supply\textquotedblright{} (Lewis \& Short); \emph{heure} is inherited from Latin \emph{hōra} \textquotedblleft{}hour\textquotedblright{}.
\end{cousinweb}

\begin{grammarlens}[title={What you've built}]
One more word for telling, explaining and describing at B1.
\end{grammarlens}

\subsection*{Your turn}
\begin{itemize}
\raggedright
  \item \emph{Say it:} \emph{les heures supplémentaires}
  \item \emph{Say it:} \emph{les heures supplémentaires}, once more
  \item \emph{Say it:} a sentence of your own with \emph{les heures supplémentaires}
  \item \emph{Recall:} say the French for the striker, then the French for the employee, then say \emph{les heures supplémentaires} again
\end{itemize}

\begin{tcolorbox}[breakable,colback=teal!4,colframe=teal!35!black,title={Before you move on}]
What does les heures supplémentaires mean? (\textquotedblleft{}The overtime\textquotedblright{}.) Say it once more.
\end{tcolorbox}
